\uuid{qFcy}
\titre{Changement de variables}
\auteur{hanani ; mzari}
\organisation{exo7}
\licence{CC BY-NC-SA 4.0}
\datecreate{2020-01-20}
\niveau{L1}
\module{Analyse}
\chapitre{Calcul d'intégrales}
\sousChapitre{Primitives diverses}
\difficulte{3}
\theme{Primitives, changement de variables}
\exo7qcmid{lille-2.primitives.044}

\contenu{
\begin{question}
\qclassification{127.11}
\qtags{motcle=changement de variables}
Pour $x\in ]-\pi/2,\pi/2[$, on pose $t=\tan (x)$ et on rappelle que $\displaystyle \sin ^2x=\frac{\tan ^2x}{1+\tan ^2x}$. Parmi les affirmations suivantes, cocher celles qui sont vraies :
\begin{answers}
\bad{$\displaystyle \frac{1}{1+\sin ^2x}=\frac{1+t^2}{1+2t^2}$ et $\displaystyle \int \frac{\mathrm{d}x}{1+\sin ^2x}=\int \frac{1+t^2}{1+2t^2}\mathrm{d}t$.}
\bad{$\displaystyle \int \frac{\mathrm{d}x}{1+\sin ^2x}=\frac{t}{2}-\frac{1}{2\sqrt{2}}\arctan(\sqrt{2}t)+k$, $k\in \Rr$.}
\good{$\displaystyle \int \frac{\mathrm{d}x}{1+\sin ^2x}=\int \frac{1}{1+2t^2}\mathrm{d}t$.}
\good{Sur $]-\pi/2,\pi/2[$, on a : $\displaystyle \int \frac{\mathrm{d}x}{1+\sin ^2x}=\frac{1}{\sqrt{2}}\arctan\left(\sqrt{2}\tan x\right)+k$, $k\in \Rr$.}
\end{answers}
\vskip2mm
\begin{explanations}
Avec $t=\tan x$, on aura $\displaystyle \frac{1}{1+\sin ^2x}=\frac{1+t^2}{1+2t^2}$ et $\displaystyle \mathrm{d}x=\frac{\mathrm{d}t}{1+t^2}$. D'où
$$\int \frac{\mathrm{d}x}{1+\sin ^2x}=\int \frac{\mathrm{d}t}{1+2t^2}=\frac{1}{\sqrt{2}}\arctan(\sqrt{2}t)+k,\; k\in \Rr.$$
\end{explanations}
\end{question}
}
